\subsection{Asymptotisches Verhalten}
\label{subsec:real_functions_asymptotes}
Nachdem wir einzelne Stellen untersucht haben, betrachten wir nun
$x\to\pm\infty$. Ein endlicher Grenzwert beschreibt bereits eine Gerade,
der sich der Graph annähert.

\par\noindent\begin{minipage}{\linewidth}
\begin{DefinitionBox}{Waagerechte Asymptote}{real_functions_asymptote}
  Existiert für $f$ der endliche Grenzwert
  \[
    \lim_{x\to\pm\infty}f(x)=L,\qquad L\in\R,
  \]
  heißt die Gerade $y=L$ eine \emph{waagerechte Asymptote}
  \EnglishTerm{horizontal asymptote} für die betreffende Richtung.
\end{DefinitionBox}
\begin{ExampleBox}{Zwei waagerechte Asymptoten}{real_functions_horizontal_asymptote_examples}
  Dividieren Sie Zähler und Nenner durch die höchste Nennerpotenz.
  Bestimmen Sie so die Asymptote für beide Richtungen im Unendlichen.
  \begin{ClozeArea}[34]{7}
    \[
      \begin{aligned}
        \lim_{x\to\pm\infty}\frac{x}{x^2+1}
          &=\lim_{x\to\pm\infty}\frac{1/x}{1+1/x^2}=0,\\[0.8em]
        \lim_{x\to\pm\infty}\frac{4x^2+2x+1}{x^2+1}
          &=\lim_{x\to\pm\infty}\frac{4+2/x+1/x^2}{1+1/x^2}=4.
      \end{aligned}
    \]
    Die waagerechten Asymptoten sind $y=0$ beziehungsweise $y=4$.
  \end{ClozeArea}
  \begin{center}
    \CachedTikzFigure{content/resources/chapters/real_functions/tikz/functions_horizontal_asymptotes.tex}
  \end{center}
\end{ExampleBox}
\end{minipage}

Ein unendlicher Grenzwert beschreibt nur das unbegrenzte Wachstum.
Eine Gerade kann den Verlauf genauer beschreiben: Wir suchen sie so,
dass der Unterschied zum Funktionsgraphen gegen null geht.

\par\noindent\begin{minipage}{\linewidth}
\begin{DefinitionBox}{Schiefe Asymptote}{real_functions_slant_asymptote_definition}
  Die Gerade $y=mx+b$ mit $m\ne0$ heißt \emph{schiefe Asymptote}
  \EnglishTerm{oblique asymptote}, wenn
  \[
    \lim_{x\to\pm\infty}\bigl(f(x)-(mx+b)\bigr)=0
  \]
  für die betreffende Richtung gilt.
\end{DefinitionBox}
\begin{ExampleBox}[before skip=0.3\baselineskip]{Die Polynomdivision zeigt die Asymptote}{real_functions_slant_asymptote}
  Wir untersuchen $f(x)=(x^2+1)/(x-1)$ für große positive und negative $x$.
  \begin{ClozeArea}[34]{6}
    Die Polynomdivision zerlegt den Term in eine Gerade und einen Rest:
    \[
      f(x)=x+1+\frac2{x-1}.
    \]
    Weil $\lim_{x\to\pm\infty}2/(x-1)=0$, ist $y=x+1$
    auf beiden Seiten eine schiefe Asymptote.
    Zusätzlich ist $x=1$ eine \textcolor{green!55!black}{vertikale Asymptote}.
  \end{ClozeArea}
  \begin{center}
    \CachedTikzFigure{content/resources/chapters/real_functions/tikz/functions_slant_asymptote.tex}
  \end{center}
\end{ExampleBox}
\end{minipage}

Bei einem größeren Gradunterschied kann sich der Graph statt einer Geraden
einem höhergradigen Polynom annähern. Die Polynomdivision liefert auch hier
den gesuchten Verlauf und einen Rest, der gegen null geht.

\begin{ExampleBox}{Annäherung an eine Parabel}{real_functions_polynomial_asymptote}
  Für $f(x)=(x^4-2x^2+1)/(x^2+1)$ ist der Zählergrad um $2$ größer.
  Bestimmen Sie durch Polynomdivision das Polynom, dem sich der Graph nähert.
  \begin{ClozeArea}[34]{9}
    \[
      \begin{array}{r@{\;}l}
        (x^4-2x^2+1):(x^2+1)&=x^2-3+\dfrac4{x^2+1}\\[0.8em]
        \underline{-(x^4+x^2)}&\\[0.5em]
        -3x^2+1&\\[0.5em]
        \underline{-(-3x^2-3)}&\\[0.5em]
        4&
      \end{array}
    \]
    Mit $A(x)=x^2-3$ gilt
    $\lim_{x\to\pm\infty}\bigl(f(x)-A(x)\bigr)=
      \lim_{x\to\pm\infty}4/(x^2+1)=0$.
    Außerdem ist $\lim_{x\to\pm\infty}f(x)=+\infty$.
  \end{ClozeArea}
  \begin{center}
    \CachedTikzFigure{content/resources/chapters/real_functions/tikz/functions_polynomial_asymptote.tex}
  \end{center}
\end{ExampleBox}

Die Beispiele lassen sich anhand der Zähler- und Nennergrade zusammenfassen.
\begin{TheoremBox}{Gradvergleich für gebrochen-rationale Funktionen}{real_functions_degree_comparison}
  Seien $m=\deg(Z)$, $n=\deg(N)$ und $a_m,b_n$ die Leitkoeffizienten.
  Für $f(x)=Z(x)/N(x)$ mit $Z\not\equiv0$ gilt:
  \begin{center}
    \renewcommand{\arraystretch}{1.5}
    \begin{tabular}{@{}ll@{}}
      \toprule
      Gradvergleich & Verhalten für $x\to\pm\infty$\\ \midrule
      $m<n$ & $\lim_{x\to\pm\infty}f(x)=0$\\
      $m=n$ & $\lim_{x\to\pm\infty}f(x)=a_m/b_n$\\
      $m=n+1$ & schiefe Asymptote aus Polynomdivision\\
      $m>n+1$ & Annäherung an ein Polynom vom Grad $m-n$\\
      \bottomrule
    \end{tabular}
  \end{center}
  Allgemein liefert die Division $Z(x)=h(x)N(x)+r(x)$ mit $\deg(r)<n$
  die Zerlegung $f(x)=h(x)+r(x)/N(x)$. Wegen $\lim_{x\to\pm\infty}r(x)/N(x)=0$ beschreibt $h$
  das Verhalten im Unendlichen.
\end{TheoremBox}
% Handschrift: PDF-Seiten 27–29 (Herleitung und sämtliche Gradvergleichsfälle).
Durch Ausklammern der höchsten Potenzen wird der Gradvergleich sichtbar:
\[
  \begin{gathered}
    f(x)=x^{m-n}
    \frac{a_m+a_{m-1}/x+\dots+a_0/x^m}
         {b_n+b_{n-1}/x+\dots+b_0/x^n},\\[0.6em]
    \lim_{x\to\pm\infty}
      \frac{a_m+a_{m-1}/x+\dots+a_0/x^m}
           {b_n+b_{n-1}/x+\dots+b_0/x^n}=\frac{a_m}{b_n}.
  \end{gathered}
\]
Für $m>n$ wächst $|f(x)|$ daher gegen $\infty$; das Vorzeichen ergibt
sich aus $a_m/b_n$ und der Parität von $m-n$.

\begin{RemarkBox}[remark title=Form und Schnittpunkte von Asymptoten]
  Für $m>n+1$ ist der Graph von $h$ eine polynomial beschriebene
  Annäherungskurve, keine schiefe Gerade. Eine Asymptote darf außerdem
  vom Funktionsgraphen geschnitten werden.
\end{RemarkBox}
